\scanpage{309}

% archive sheet number: 308

\[
\tau_1\xrightarrow{f_{\deleted{\unsure{1}}}}\tau_2,\quad
\tau'_1\xrightarrow{f'_{\deleted{\unsure{2}}}}\tau'_2
\qquad\text{donne}\qquad
\tau_1\tau'_1\xrightarrow{f_{\deleted{\unsure{1}}}*f'_{\deleted{\unsure{2}}}}\tau_2\tau'_2
\]
(on \unsure{n'utilisera|utilisera} pas $f_{\deleted{\unsure{1}}}\overset{\unsure{*}}{\unsure{\cdot|\circ}}f'_{\deleted{\unsure{2}}}$ pour la flèche notée $f_{\deleted{\unsure{1}}}*f'_{\deleted{\unsure{2}}}$, à cause de la confusion possible avec la composition des flèches dans $\underline{\mathrm{Ty}}$) : On a transitivité de façon évidente si de plus $\tau_2\overset{g_{\deleted{\unsure{1}}}}{\unsure{\text{---}|\longrightarrow}}\tau_3$, $\tau'_2\overset{g'}{\unsure{\text{---}|\longrightarrow}}\tau'_3$
\[
\begin{aligned}
&\deleted{\unsure{(g_1f_1)*(g_2|(gf)*(g'|(gf)*g'}}\\
&(gf)*(g'f')=(g*g')\circ(f*f').
\end{aligned}
\]

Donc on a un bifoncteur de composition
\[
\underline{\mathrm{Ty}}\times\underline{\mathrm{Ty}}\longrightarrow\underline{\mathrm{Ty}}
\]
\margin{$(c,c')\longmapsto c.c'$\\ ou $c*c'$}

(En fait, \deleted{\illegible{2}}\inserted{on a aussi} un bifoncteur \deleted{$\unsure{\mathrm{Ty}|\mathrm{Typ}}$}
\[
\underline{\mathrm{Typ}}\times\underline{\mathrm{Typ}}\unsure{\text{ --- }|\longrightarrow}\underline{\mathrm{Typ}},
\]
mais on préfère oublier maintenant la catégorie $\underline{\mathrm{Typ}}$, qui donne un formalisme nettement moins simple.)

Cette composition est \uline{strictement} associative et \uline{strictement} \uline{unitaire} (l'unité étant bien sûr l'élément unité $\unsure{1|s}$ du monoïde libre $\mathrm{Ty}$, correspondant à la catégorie ponctuelle $\{0\}$), et pas seulement ``à isom près''. Et pour cause, puisque dans $\underline{\mathrm{Ty}}$ les \deleted{\unsure{objets}} iso sont des identités. \uline{Donc} $\underline{\mathrm{Ty}}$ \uline{est un monoïde dans la catégorie} $(\mathrm{Cat})$ des petites catégories!
